% (C) 2026 Ralf Stephan, in collaboration with Claude Code.
% Released under CC0 1.0 Universal (public-domain dedication).
% See https://creativecommons.org/publicdomain/zero/1.0/

\documentclass[11pt]{article}
\usepackage[T1]{fontenc}
\usepackage[utf8]{inputenc}
\usepackage{lmodern}
\usepackage[margin=2.5cm]{geometry}
\usepackage{amsmath,amssymb}
\usepackage{booktabs,tabularx}
\usepackage{braket}
\usepackage[hidelinks]{hyperref}

% Lean names and paths: verbatim, breakable at . _ /
\DeclareUrlCommand\lean{\urlstyle{tt}}
\def\UrlBreaks{\do\.\do\_\do\/}
\newcolumntype{L}{>{\raggedright\arraybackslash}X}
\newcolumntype{B}{>{\raggedright\arraybackslash}p{0.42\textwidth}}
\emergencystretch=6em

\title{Lean formalization of an Assumptions of Physics technical brief:\\ log}
\author{}
\date{}

\begin{document}
\maketitle

\noindent
This log records which statements of the technical brief in this directory were checked in Lean,
where the formal version differs from the brief, and any problems found in the brief.
The brief is available at \url{https://assumptionsofphysics.org/resources/briefs}.

\section*{Setup}

\begin{itemize}
\item Lean \texttt{v4.36.0-rc1}; Mathlib at commit \texttt{2cabef97}; physlib (\texttt{master})
  and TauCeti (commit \texttt{eb4dd24}, on the same Mathlib) as further dependencies.
\item Build: \texttt{lake build} in the project root.
\item Every file is checked with \texttt{\#print axioms} on its main results. A result is
  recorded here only if it depends on no axioms beyond \lean{propext},
  \lean{Classical.choice} and \lean{Quot.sound}, and so has no \lean{sorry}.
\item Imports are kept to the Mathlib modules actually needed.
\item Material that belongs in Mathlib is kept in \lean{ForMathlib/}, in Mathlib's directory
  layout (library \lean{ForMathlib}), and imported by the brief files.
\end{itemize}


%======================================================================

\section{Non-linear representations in quantum mechanics}

\paragraph{Files.}\leavevmode
\begin{itemize}
\item File: \lean{briefs/NonLinearQuantum/NonLinearQuantum.lean}
\item Library: \lean{NonLinearQuantum}
\item Imports: \lean{Mathlib.Analysis.InnerProductSpace.Calculus},\newline
  \lean{Mathlib.Analysis.SpecialFunctions.ExpDeriv}
\end{itemize}

\paragraph{Setting.} $H = \mathbb{C}^2$ (\texttt{EuclideanSpace} $\mathbb{C}$ \texttt{(Fin 2)}), with
$\ket{z^+} = e_0$, $\ket{z^-} = e_1$, and
\[
  m(v) = e^{i\theta\, p(z^+\mid v)}\, v, \qquad p(z^+\mid v) = \frac{|v_0|^2}{\|v\|^2}.
\]
In Lean $m(0) = 0$ holds without a case split, since $0/0 = 0$.

\paragraph{Results.}
\begin{center}
\begin{tabularx}{\textwidth}{@{}BL@{}}
\toprule
Brief & Lean \\
\midrule
eq.~(3), rays to rays & \lean{nlMap_smul}: $m(kv) = k\,m(v)$ for every $k\in\mathbb{C}$ \\
each ray is mapped to itself & \lean{nlMap_mem_ray} \\
eq.~(4), $m$ not linear & \lean{nlMap_not_add}, \lean{not_isLinearMap_nlMap}, for $\theta\in(0,2\pi)$ \\
eq.~(5), inner product & \lean{inner_nlMap} \\
eq.~(6), Born rule preserved & \lean{born_nlMap} \\
Remark: $m$ bijective & \lean{nlHomeomorph}: a homeomorphism of $\mathbb{C}^2$, inverse $m_{-\theta}$ \\
Remark: $m$ a diffeomorphism & \lean{contDiffAt_nlMap}: smooth at every $v \ne 0$ \\
(correction) & \lean{not_differentiableAt_nlMap_zero}: not differentiable at $0$ \\
Theorem~1, mathematical content & \lean{exists_nonlinear_born_preserving} \\
Born rule depends only on rays & \lean{born_smul}: $p(cv\mid dw) = p(v\mid w)$ for $c, d \ne 0$ \\
$m(H)$, structure carried along $m$ & \lean{Twisted}, \lean{toTwisted} (identity of the carrier) \\
same scalars, hence same rays & \lean{toTwisted_smul} \\
same Born rule & \lean{born_toTwisted} \\
different addition & \lean{not_map_add_toTwisted} \\
(correction) $H \cong m(H)$ & \lean{twistedIsometry}: $m$ is a linear isometry onto $m(H)$ \\
Theorem~1, precise form & \lean{exists_second_linear_structure} \\
(caveat) dimension one & \lean{map_add_of_map_smul}: maps commuting with scalars are additive \\
\bottomrule
\end{tabularx}
\end{center}

\paragraph{Differences from the brief.} Theorem~1 (``REP-VSLIN is independent of QST'') is a
statement about the logical independence of two conditions, which the brief does not make
precise. Two things are checked.
\begin{itemize}
\item The existence of a homeomorphism of $H$ that maps every ray to itself, preserves the Born
  rule and is not additive.
\item A precise form of $\neg$REP-VSLIN under which the brief's proof is a proof: the linear
  structure of $H$ is not fixed by the rays and the Born rule. The type \lean{Twisted}~$\theta$
  is the carrier of $H$ with addition $v \oplus w = m(m^{-1}v + m^{-1}w)$, scalar multiplication
  $k \odot v = m(k\,m^{-1}v)$ and inner product $\langle m^{-1}v, m^{-1}w\rangle$; it is a complex
  inner product space. By eq.~(3) its scalar multiplication is that of $H$, so it has the same
  rays; by eq.~(6) it has the same Born rule; by eq.~(4) its addition differs from that of $H$.
\end{itemize}
See the open problems for why the brief has to choose this reading.

\paragraph{Problems found in the brief.}
\begin{itemize}
\item The closing remark says that $m$ is a diffeomorphism. It is not differentiable at $0$:
  along every real line $m(tv) = t\,m(v)$, so a derivative at $0$ would have to equal $m$
  itself, and $m$ is not additive. The correct statement is that $m$ is a homeomorphism of
  $\mathbb{C}^2$ and a diffeomorphism of $\mathbb{C}^2\setminus\{0\}$. Theorem~1 only needs a
  bijection, so it is not affected.
\item The proof says that ``since the map is non-linear, $H \not\equiv m(H)$ in the category of
  vector spaces''. This is false. As a set $m(H) = H$, since $m$ is a bijection. With the vector
  space structure of $H$, $m(H)$ is $H$. With the structure carried over along $m$, $m$ is by
  construction a linear isometry $H \to m(H)$ (\lean{twistedIsometry}). What is true is that
  the identity of the carrier is not linear from one structure to the other.
\item Eq.~(6) holds for every map that sends each vector into its own ray, since the Born rule
  only sees rays (\lean{born_smul}). Any $v \mapsto \varphi([v])\,v$ with
  $\varphi \colon P(H) \to \mathbb{C}\setminus\{0\}$ would do. What matters for the proof is
  that $m$ commutes with scalars (eq.~(3)) and is not additive (eq.~(4)).
\item QST allows $\dim H \le 1$. There every map that commutes with scalars is additive
  (\lean{map_add_of_map_smul}), so there is no second linear structure. Theorem~1 claims only
  consistency, so the two-state example suffices; ``for every system'' would need
  $\dim H \ge 2$.
\end{itemize}

\end{document}
